\documentclass[main.tex]{subfiles}
\begin{document}
\ifSubfilesClassLoaded{\begin{center}\small\itshape Part of the GnssProcessor Mathematical Reference \textbullet\ \today\end{center}}{}
\section{Observation Model Physics: Antenna, Attitude, Relativity, and Range Corrections}\label{sec:obsphys}
% ============================================================
%  Deterministic per-observation corrections beyond geometry
%  and atmosphere: antenna PCO/PCV, satellite attitude, carrier
%  wind-up, relativity, Shapiro, Sagnac, and group-delay/bias
%  realisation of the abstract B/b terms.
% ============================================================

% doc-local macros (defined in the body, never in a section title)
\newcommand{\bframe}{\bm{F}}                       % satellite body-frame matrix [i_s j_s k_s]
\newcommand{\isun}{\hat{\bm{x}}_\sat}              % sat body +X (toward Sun)
\newcommand{\jcross}{\hat{\bm{y}}_\sat}            % sat body +Y (cross-track)
\newcommand{\knad}{\hat{\bm{k}}_\sat}              % sat body +Z (nadir / boresight)
\newcommand{\esun}{\hat{\bm{e}}_\odot}            % unit Sun direction
\newcommand{\pco}{\bm{p}}                          % phase-centre offset vector
\newcommand{\dipS}{\bm{D}_\sat}                    % satellite effective dipole
\newcommand{\dipR}{\bm{D}_\rx}                     % receiver effective dipole
\newcommand{\winddphi}{\Delta\varphi}             % signed wind-up angle (rad)

This document covers the \emph{deterministic} corrections of the
observation model, other than geometry, the atmosphere, and the streamed
SSR overlay.  \dcref{sec:obs}{The Foundations observation equations}
(\dcref{eq:obs-code}{the undifferenced code equation} and
\dcref{eq:obs-phase}{the undifferenced phase equation}) introduce the
satellite and receiver hardware biases $B^\sat_\sigl,\,b^\sat_\sigl,
\,B_{\rx,\sigl},\,b_{\rx,\sigl}$ only as symbols.
\dcref{sec:obs-geometry}{The observation-geometry section} gives the
geometric range $\geomrange^\sat_\rx$ and its linearisation.  This
document supplies every other deterministic term between the raw
measurement and the modelled range:
\begin{itemize}
\item the antenna phase-centre geometry, and the satellite attitude
      frame it depends on;
\item the carrier-phase wind-up;
\item the relativistic clock and propagation corrections;
\item the Sagnac (Earth-rotation) frame correction;
\item the bias terms, realised through broadcast group delays and the
      precise \texttt{.BIA} product.
\end{itemize}
The streamed SSR overlay corrects the same terms in real time and is
described in \dcref{sec:ssr}{the SSR corrections document}.  The two
mechanisms differ in frames, sign conventions, and gating, and the
differences are noted where they matter.

\texttt{PppModel::computeBaseModeledValue} in
\texttt{model/PppModel.cpp} assembles every term below into one modelled
observable.  Both the iono-free-pair and the uncombined observation
paths call it.  It adds the terms to the geometric range from
\texttt{model/ModelGeometry.cpp}.  For the code observable the sum is
\begin{equation}
\hat{\code}^\sat_{\rx,\sigl}
=
\geomrange^\sat_\rx
\;-\;\sLight\,\mathrm{d}t^\sat
\;-\;\sLight\,\delta t_{\mathrm{rel}}
\;+\;T^\sat_\rx
\;+\;A^\sat_{\rx,\sigl}
\;+\;B_{\mathrm{grp}}
\;+\;S^\sat_\rx
\;+\;d_{\mathrm{atm}}
\;+\;d_{\mathrm{tide}},
\label{eq:obsphys-assembly}
\end{equation}
where $A^\sat_{\rx,\sigl}$ is the total antenna correction
(\Cref{sec:obsphys-antenna}), $\delta t_{\mathrm{rel}}$ the periodic
relativistic clock term (\Cref{sec:obsphys-rel}), $S^\sat_\rx$ the
Shapiro delay (\Cref{sec:obsphys-shapiro}), $B_{\mathrm{grp}}$ the
deterministic code/phase bias (\Cref{sec:obsphys-bias}), and
$d_{\mathrm{atm}},\,d_{\mathrm{tide}}$ the loading/site-displacement
terms documented in \dcref{sec:sitedisp}{the site-displacement
document}.  The Sagnac correction (\Cref{sec:obsphys-sagnac}) is part of
$\geomrange^\sat_\rx$ and is not added separately; the reason is given
below.  The signs in \eqref{eq:obsphys-assembly} follow the return
statement of \texttt{PppModel::computeBaseModeledValue}
(\texttt{model/PppModel.cpp}): the satellite clock and relativistic terms
are \emph{subtracted}, and the antenna, Shapiro, loading and tide terms
are \emph{added}.

% ============================================================
\subsection{Satellite attitude: the body frame}\label{sec:obsphys-attitude}
% ============================================================

The satellite phase-centre offset, the phase-centre variation, and the
carrier wind-up are all defined in the \emph{spacecraft body frame}, so
that frame is constructed first.
The engine uses the IGS/ANTEX convention: a right-handed orthonormal
triad $(\isun,\jcross,\knad)$ with $+Z$ toward the Earth's centre
(nadir / antenna boresight), $+X$ in the Sun--satellite--Earth plane
pointing toward the Sun, and $+Y$ completing the triad along the solar-panel
(cross-track) axis (\texttt{model/SatAttitude.cpp},
\texttt{model/Antenna.cpp}).

\paragraph{Nominal yaw-steering.}
With $\satpos$ the satellite ECEF position and
$\esun=(\bm{r}_\odot-\satpos)/\lVert\bm{r}_\odot-\satpos\rVert$ the unit
direction to the Sun, the nadir axis and the cross-track axis are
\begin{equation}
\knad = -\,\frac{\satpos}{\lVert\satpos\rVert},
\qquad
\jcross = \frac{\knad\times\esun}{\lVert\knad\times\esun\rVert},
\qquad
\isun = \jcross\times\knad,
\label{eq:obsphys-yaw-frame}
\end{equation}
in the operand order of \texttt{SatAttitude::compute}
(\texttt{model/SatAttitude.cpp}).  The antenna path calls that routine
(\texttt{Antenna::getAntennaCorrection}) and does not rebuild the frame.
The wind-up path does contain an inline copy of the construction, noted
in \Cref{sec:obsphys-windup}.  Forming $\jcross$ first and then
$\isun=\jcross\times\knad$ is the Gram--Schmidt orthogonalisation of the
Sun direction against nadir.  By the
$\bm{a}\times(\bm{b}\times\bm{c})$ identity,
\begin{equation}
\isun = \jcross\times\knad
      = (\knad\times\esun)\times\knad
      = \esun - \knad\,(\knad\cdot\esun),
\label{eq:obsphys-xsun-gramschmidt}
\end{equation}
so $\isun$ is the normalised component of $\esun$ orthogonal to nadir.
The $+X$ axis therefore points toward the Sun as far as the nadir
constraint allows, as the ANTEX convention requires.  The body-frame
matrix with these three columns,
\begin{equation}
\bframe = \big[\,\isun\;\;\jcross\;\;\knad\,\big],
\label{eq:obsphys-bodyframe-matrix}
\end{equation}
is orthonormal, so its inverse is its transpose.  This is used below to
rotate a body-frame offset into ECEF.

\paragraph{Sun elevation and the mode switch.}
Non-nominal attitude (eclipse seasons, the noon and midnight yaw turns)
depends on the Sun's elevation above the orbital plane, $\beta$.  With the orbit angular-momentum (normal) unit vector
$\hat{\bm{n}}=(\satpos\times\satvel)/\lVert\satpos\times\satvel\rVert$,
\begin{empheq}[box=\fbox]{equation}
\beta = \arcsin\!\big(\hat{\bm{n}}\cdot\esun\big),
\qquad
\hat{\bm{n}} = \frac{\satpos\times\satvel}{\lVert\satpos\times\satvel\rVert},
\label{eq:obsphys-beta}
\end{empheq}
signed by the projection of the Sun direction onto the $+(\satpos\times
\satvel)$ normal (\texttt{SatAttitude::computeBeta},
\texttt{model/SatAttitude.cpp}; the sign is noted in-code as unvalidated
against an external attitude product).  The mode selector
(\texttt{SatAttitude::determineMode}) keeps GPS, GLONASS and Galileo in
yaw-steering unconditionally, and switches BeiDou and QZSS to
orbit-normal attitude below a constellation-specific $\beta$ threshold
($4.0^\circ$ BeiDou, $20.0^\circ$ QZSS,
\texttt{model/SatAttitude.cpp}).  In orbit-normal mode the $+Y$ axis is
pinned to the orbit normal, $\jcross=\hat{\bm{n}}$, and the triad is
re-orthogonalised $\isun=\jcross\times\knad$, $\jcross=\knad\times\isun$.
The legacy default path also forces a hard-coded list of BeiDou-2 GEO
PRNs (C01--C05, C59, C60) into orbit-normal mode whatever the value of
$\beta$.  The $\beta$-driven table is active only when
\texttt{enableOrbitNormalAttitude()} is set.  With the flag off the
result is bit-identical to the earlier inline frame
(\texttt{model/Antenna.cpp}, \texttt{model/SatAttitude.cpp}).

\paragraph{Eclipse gating.}
The body frame above assumes a known yaw angle.  That assumption fails
during Earth-shadow passes.  The actual yaw during eclipse and the
recovery after it depend on the satellite block and are not modelled
here, so the antenna PCO/PCV and the wind-up are unreliable for an
eclipsing satellite.  When the user opts in
(\texttt{processing.exclude\_eclipsing\_satellites}), the satellite is
dropped at the geometry stage by a \emph{cylindrical umbra} test
(\texttt{Planetary::isSatelliteEclipsing}, \texttt{model/Planetary.cpp},
called from \texttt{model/ModelGeometry.cpp}).  With
$\hat{\bm{u}}_\odot=\bm{r}_\odot/\lVert\bm{r}_\odot\rVert$ the Earth-to-Sun
direction, the satellite's along-axis and perpendicular distances are
$s_\parallel=\satpos\cdot\hat{\bm{u}}_\odot$ and $\bm{s}_\perp=\satpos -
s_\parallel\,\hat{\bm{u}}_\odot$, and the predicate is
\begin{equation}
\text{eclipsed}
\;\Longleftrightarrow\;
\big(s_\parallel < 0\big)
\;\wedge\;
\big(\lVert\bm{s}_\perp\rVert < \kappa\,R_\oplus\big),
\label{eq:obsphys-eclipse}
\end{equation}
i.e. the satellite is on the anti-Sun side \emph{and} inside the shadow
cylinder of radius $\kappa R_\oplus$, with $R_\oplus = 6\,371\,230$\,m
(\texttt{common/Constants.hpp}) and $\kappa$ the user shadow scale
(\texttt{processing.eclipse\_shadow\_scale}, default unity; cited rather
than reproduced as a tuning knob).  This is a conservative geometric
umbra test.  It does not model the penumbra.

% ============================================================
\subsection{Antenna phase-centre corrections}\label{sec:obsphys-antenna}
% ============================================================

The range is measured between the satellite and receiver \emph{antenna
phase centres}.  The estimated states refer to the satellite centre of
mass and the receiver antenna reference point (ARP).  The position state
must be the ARP because the ANTEX receiver PCO/PCV are referenced to it.
The antenna correction $A^\sat_{\rx,\sigl}$ connects the two, as the sum
of four geometric terms and the wind-up:
\begin{empheq}[box=\fbox]{equation}
A^\sat_{\rx,\sigl}
\;=\;
-\,\delta_{\mathrm{PCO}}^{\sat}
\;+\;\delta_{\mathrm{PCV}}^{\sat}
\;+\;\delta_{\mathrm{PCO}}^{\rx}
\;+\;\delta_{\mathrm{PCV}}^{\rx}
\;+\;\varphi_{\mathrm{wu}},
\label{eq:obsphys-antenna-total}
\end{empheq}
as accumulated in \texttt{Antenna::getAntennaCorrection}
(\texttt{model/Antenna.cpp}).  The line of sight used in this section is
the \emph{satellite-to-receiver} unit vector,
$\los = (\rxpos_\rx-\satpos)/\lVert\rxpos_\rx-\satpos\rVert$
(\texttt{Antenna.cpp}).  It is the negative of the receiver-to-satellite
$\los^\sat_{\hat{\rx}}$ of \dcref{sec:obs-geometry}{the geometry section}
(\dcref{eq:los}{its line-of-sight definition}).  The sign of each term
below follows from this choice.

\paragraph{Which coordinate is the ARP.}
Every position the engine derives from a RINEX \emph{header} is an ARP
position.  The marker coordinate is reduced to the ARP through the
header eccentricity before it seeds the state
(\texttt{Coordinates::markerToArp}, packaged as
\texttt{RinexParser::getArpPositionXyz}; \dcref{sec:frame-arp}{the
marker/ARP section of the Coordinate Frames document}).  A coordinate
supplied directly in the configuration is \emph{not} reduced.  It is
used as given, on the policy that the user owns its reference point.
This matters in the network producer.  Its station coordinates come from
the configuration (\texttt{StationConfig::position\_ecef}) and are held
fixed, with their position columns dropped from the normal equations.  A
monument coordinate there is not absorbed by a position state.  The
offset goes into the residuals, and from there into the estimated
clocks, biases, and ionosphere.

\paragraph{Satellite PCO (subtracted).}
The satellite phase-centre offset $\pco^\sat_{\mathrm{body}}$ is published
in the body frame (millimetres in ANTEX).  It is rotated to ECEF by the
body-frame matrix; because $\bframe$ collects the body axes as
\emph{columns}, $\bframe$ itself maps body $\to$ ECEF (and the orthonormal
$\bframe^\Tr=\bframe^{-1}$ maps ECEF $\to$ body).  It is then projected onto the
LOS:
\begin{equation}
\delta_{\mathrm{PCO}}^{\sat}
=
\big(\bframe\,\pco^\sat_{\mathrm{body}}\big)\cdot\los .
\label{eq:obsphys-sat-pco}
\end{equation}
The antenna phase centre is \emph{nearer the receiver} than the centre
of mass, so the true range is shorter than the CoM range and the
correction is \emph{subtracted} in \eqref{eq:obsphys-antenna-total}
(\texttt{correction -= sat\_pco\_corr}, \texttt{model/Antenna.cpp}).
There is one real-time exception.  When an applied SSR orbit correction
is referenced to the antenna phase centre, the corrected orbit already
points at the phase centre.  The \texttt{skip\_satellite\_pco} path of
\texttt{getAntennaCorrection} then drops the satellite-PCO term so that
it is not counted twice.  The path is gated by
\texttt{realtime.honor\_ssr\_reference\_point} and is off by default.
The SSR reference-point convention is documented in
\dcref{sec:ssr}{the SSR corrections document}.

\paragraph{Satellite PCV (added).}
The satellite phase-centre variation is a function of the \emph{nadir
angle}, the angle between the boresight $\knad$ and the line of sight,
and of the azimuth in the body frame.  With $\cos\theta_{\mathrm{nad}} =
\knad\cdot\los$,
\begin{equation}
\theta_{\mathrm{nad}} = \arccos\!\big(\knad\cdot\los\big),
\qquad
\alpha_{\mathrm{body}} = \operatorname{atan2}\!\big((\bframe^\Tr\,\los)_y,\,
(\bframe^\Tr\,\los)_x\big),
\label{eq:obsphys-nadir}
\end{equation}
and the PCV value is read by bilinear interpolation of the ANTEX grid
over (zenith/nadir ring, azimuth sector), the look-up argument being
$\theta_{\mathrm{nad}}-\mathtt{zen1}$ in grid units
(\texttt{Antenna::interpolatePcv}, \texttt{model/Antenna.cpp}).  The
satellite PCV is \emph{added} to the range correction.

\paragraph{Receiver PCO (added).}
The receiver PCO is published as a North/East/Up (NEU) triplet relative to
the ARP.  It is reordered to ENU, rotated to ECEF by the local
ENU$\to$ECEF rotation $\bm{R}_{\mathrm{enu}}(\phi,\lambda)$ at the
receiver geodetic latitude/longitude
(\texttt{Coordinates::enuToEcefMatrix}), and projected on the LOS:
\begin{equation}
\delta_{\mathrm{PCO}}^{\rx}
=
\Big(\bm{R}_{\mathrm{enu}}(\phi,\lambda)\,
\big[\,p_E\;\;p_N\;\;p_U\,\big]^\Tr\Big)\cdot\los,
\qquad
\big[\,p_E\;p_N\;p_U\,\big] = \big[\,p_{\mathrm{NEU},2}\;\;
p_{\mathrm{NEU},1}\;\;p_{\mathrm{NEU},3}\,\big],
\label{eq:obsphys-rx-pco}
\end{equation}
where the index swap $(N,E,U)\!\to\!(E,N,U)$ is explicit in
\texttt{model/Antenna.cpp}.  It is \emph{added} to the correction.

\paragraph{Receiver PCV (added).}
The receiver PCV is indexed by \emph{zenith angle} $z = 90^\circ -
\mathrm{el}$ (and, optionally, azimuth corrected for the antenna mounting
orientation when \texttt{enableRxAntennaOrientation()} is set), with the
same bilinear interpolation over $z-\mathtt{zen1}$
(\texttt{model/Antenna.cpp}).  It is \emph{added} to the correction.

\paragraph{ATX frequency fallback.}
ANTEX entries often omit the newest carriers: L5/E5a/B2a on older
spacecraft, and many receiver antennas published before triple-frequency
tracking.  Setting the antenna correction to zero would leave a vertical
bias of several centimetres.  The engine instead substitutes the
\emph{nearest carrier} when the requested ANTEX frequency key is absent
(\texttt{pickAtxFrequencyKey},
\texttt{model/Antenna.cpp}).  The chains are ordered by carrier-frequency
proximity, e.g. for GPS $\mathtt{G05}\to\{\mathtt{G02},\mathtt{G01}\}$ and
for Galileo $\mathtt{E05}\to\{\mathtt{E08},\mathtt{E07},\mathtt{E06},
\mathtt{E01}\}$.  The PCO of one physical antenna typically changes by
less than a centimetre between frequencies.  The PCO and PCV chains are
searched \emph{independently}, because an antenna may have a PCO but no
PCV at a given key.

For the \emph{satellite} antenna only the chain within the same system
is allowed; another constellation's spacecraft entry would be
meaningless.  For the \emph{receiver} antenna a \emph{cross-system}
chain to the nearest GPS band is allowed after the same-system chain is
exhausted, and only when
\texttt{enable\_receiver\_cross\_system\_atx\_fallback} is set
(\texttt{pickReceiverAtxFrequencyKey}).  This is justified because the
receiver phase pattern depends on frequency, not on constellation.  The
important mappings are exact carrier matches
($\mathtt{E01}\!\to\!\mathtt{G01}$, $\mathtt{E05}\!\to\!\mathtt{G05}$ at
$1575.42$ and $1176.45$\,MHz).  The rest are approximations over
$\sim$10--70\,MHz (\texttt{model/Antenna.cpp}).  Each substitution sets a
coverage flag and writes a log line.  The same fallback policy is used
for biases: substitute and report, and do not exclude the measurement.

% ============================================================
\subsection{Carrier-phase wind-up}\label{sec:obsphys-windup}
% ============================================================

A right-hand circularly polarised GNSS signal gains a phase shift equal
to the relative rotation of the transmitting and receiving antenna
dipoles about the line of sight \cite{wuebbena1985software}.  This
wind-up is the same in cycles on every signal, so as a range it is
proportional to wavelength.  It is computed in
\texttt{GnssCorrections::computePhaseWindUp}
(\texttt{corrections/GnssCorrections.cpp}).  With $\bm{k}$ the unit
line-of-sight, $(\isun,\jcross)$ the satellite body $\pm X,+Y$ axes and
$(\bm{x}_\rx,\bm{y}_\rx)$ the receiver local North/East axes, the
\emph{effective dipoles} are
\begin{equation}
\dipS = \isun - \bm{k}\,(\bm{k}\cdot\isun) - \bm{k}\times\jcross,
\qquad
\dipR = \bm{x}_\rx - \bm{k}\,(\bm{k}\cdot\bm{x}_\rx) - \bm{k}\times\bm{y}_\rx,
\label{eq:obsphys-windup-dipoles}
\end{equation}
which is the standard Wu construction.  Each dipole is the antenna $X$
axis projected into the plane orthogonal to $\bm{k}$, \emph{minus} the
$\bm{k}\times Y$ quadrature term (\texttt{corrections/GnssCorrections.cpp}).

The sign of the quadrature term is fixed by the antenna axes.  It is
$-\operatorname{sgn}\!\big((\bm{x}\times\bm{y})\cdot\bm{k}\big)$, which is
negative when the antenna's third axis points along $+\bm{k}$.  That is
the case for both triads here: the satellite's $\bm{e}_z$ is the
geocentric nadir, and the receiver's implied third axis is
$\mathrm{North}\times\mathrm{East}=\mathrm{Down}$.  The wrong sign gives
$|\bm{D}| = 1-\cos\varepsilon$, a dipole that \emph{vanishes on
boresight}.  RTKLIB's $(-,+)$ pairing is stated with
$\bm{y}_\rx = -\mathrm{East}$ and is the same construction; with the
$+\mathrm{East}$ convention used here it becomes $(-,-)$.
\emph{Earlier revisions of this document printed $(+,+)$, copied from a
defect in the code that was corrected on 2026-08-21; see
\texttt{docs/SPEC\_phase\_windup\_sign.md}.}

The receiver North and East axes are built from the geodetic latitude
$\phi$ and longitude $\lambda$,
$\bm{x}_\rx = (-\sin\phi\cos\lambda,\,-\sin\phi\sin\lambda,\,\cos\phi)$ and
$\bm{y}_\rx = (-\sin\lambda,\,\cos\lambda,\,0)$
(\texttt{GnssCorrections.cpp}).  The signed wind-up angle is the
$\operatorname{atan2}$ of the cross and dot products of the unit
dipoles.  This covers the full range $(-\pi,\pi]$ and avoids the loss of
precision of an $\arccos$ near $0$ and $\pi$:
\begin{empheq}[box=\fbox]{equation}
\winddphi
=
\operatorname{atan2}\!\Big(
\bm{k}\cdot\big(\hat{\dipS}\times\hat{\dipR}\big),\;
\hat{\dipS}\cdot\hat{\dipR}
\Big),
\qquad
\varphi_{\mathrm{wu}}
=
\winddphi\,\frac{\lambda_\sigl}{2\pi},
\label{eq:obsphys-windup-angle}
\end{empheq}
with $\hat{\dipS},\hat{\dipR}$ the normalised dipoles and $\lambda_\sigl$
the signal wavelength (\texttt{corrections/GnssCorrections.cpp}); the
$\bm{k}\!\cdot\!(\hat{\dipS}\times\hat{\dipR})$ argument fixes the sign of
the rotation about the line of sight.  The result enters
\eqref{eq:obsphys-antenna-total} as $\varphi_{\mathrm{wu}}$ on phase
observables only (the code branch passes zero wavelength).

\paragraph{Fractional vs.\ accumulated modes.}
The $\operatorname{atan2}$ returns the wind-up wrapped to one cycle.  The
default \emph{fractional} mode uses this wrapped angle directly.  The
float ambiguity $N^\sat_{\rx,\sigl}$ absorbs any whole-cycle part, so
the wrapping does no harm.  Integer ambiguity resolution needs the
ambiguity to stay a constant integer across the arc.  For that the
engine offers an \emph{accumulated} (unwrapped) mode
(\texttt{config.accumulatePhaseWindUp()}), which unwraps each epoch
relative to the previous stored angle,
$\winddphi_t = \winddphi_{t-1} + \operatorname{wrap}_{2\pi}(\winddphi_t -
\winddphi_{t-1})$, provided the gap between epochs is below the $300$\,s
continuity threshold.  A larger gap, or a backward time step, restarts
the accumulator from the raw angle
(\texttt{corrections/GnssCorrections.cpp}).  A gap that long invalidates
the unwrap reference, so the engine starts a new arc, as it does for an
ambiguity.

\paragraph{Known frame duplication.}
The satellite body frame feeding the wind-up dipole $\dipS$
\eqref{eq:obsphys-windup-dipoles} is built \emph{inline} inside
\texttt{computePhaseWindUp} (\texttt{corrections/GnssCorrections.cpp})
rather than calling the shared \texttt{model/SatAttitude.cpp} helper used
by the antenna path \eqref{eq:obsphys-yaw-frame}.  The two constructions
are kept bit-for-bit identical: the same operand order, the same
BeiDou-2 GEO orbit-normal special case, and the same velocity fallback.
Both source files note the duplication.  It is recorded here as a
\emph{maintenance} obligation.  A change to one frame must be copied to
the other.  Otherwise the antenna PCV and the wind-up would use
different frames for the same satellite, with no error reported.

% ============================================================
\subsection{Periodic relativistic clock correction}\label{sec:obsphys-rel}
% ============================================================

A clock on an eccentric orbit runs at a varying rate relative to a clock
on the geoid.  The broadcast or precise clock polynomial absorbs the
constant and secular parts.  The \emph{periodic} part, caused by orbital
eccentricity, must be applied separately.  The standard expression
\cite{subirana2013gnss}, in
\texttt{GnssCorrections::computeRelativisticCorrection}
(\texttt{corrections/GnssCorrections.cpp}), is
\begin{empheq}[box=\fbox]{equation}
\delta t_{\mathrm{rel}}
\;=\;
-\,\frac{2\,(\satpos\cdot\satvel)}{\sLight^{2}}
\quad[\text{s}],
\qquad
\text{range term } = -\,\sLight\,\delta t_{\mathrm{rel}},
\label{eq:obsphys-rel}
\end{empheq}
applied to the range by the $-\sLight\,\delta t_{\mathrm{rel}}$ term of
the assembly \eqref{eq:obsphys-assembly} (\texttt{model/PppModel.cpp}).
Two points need stating.

First, the dot product $\satpos\cdot\satvel$ is the same whether the
position and velocity are expressed in the inertial frame or in the
rotating ECEF frame.  The ECEF velocity is
$\satvel_{\mathrm{ecef}} = \satvel_{\mathrm{inertial}} -
\bm{\omega}_\oplus\times\satpos$, and
$\satpos\cdot(\bm{\omega}_\oplus\times\satpos)=0$ because the cross
product is orthogonal to $\satpos$.  The ECEF velocity from the
broadcast ephemeris can therefore be used directly.

Second, the periodic term is computed here because the IGS precise CLK
product, by convention, contains only the constant and drift parts of
the clock.  Applying \eqref{eq:obsphys-rel} supplies the periodic part
consistently for both precise and broadcast clock sources.

% ============================================================
\subsection{Shapiro (gravitational-delay) correction}\label{sec:obsphys-shapiro}
% ============================================================

The Earth's gravitational potential slows and bends the signal.  The
resulting Shapiro delay is a few centimetres
\cite{subirana2013gnss}.  With $r_\sat=\lVert\satpos\rVert$,
$r_\rx=\lVert\rxpos_\rx\rVert$ the geocentric radii and
$r_{\sat\rx}=\lVert\satpos-\rxpos_\rx\rVert$ the satellite--receiver
distance, the delay (in metres) is
\begin{empheq}[box=\fbox]{equation}
S^\sat_\rx
\;=\;
\frac{2\,GM}{\sLight^{2}}\,
\ln\!\frac{r_\sat + r_\rx + r_{\sat\rx}}
          {r_\sat + r_\rx - r_{\sat\rx}},
\label{eq:obsphys-shapiro}
\end{empheq}
with $GM = 3.986005\times10^{14}\,\mathrm{m^3/s^2}$
(\texttt{Constants::GRAVITATIONAL\_CONSTANT}, \texttt{common/Constants.hpp})
and the formula as coded in \texttt{GnssCorrections.cpp}.  Since
$r_{\sat\rx}<r_\sat+r_\rx$ the logarithm is positive, so $S^\sat_\rx>0$
and the term is \emph{added} to the modelled range in
\eqref{eq:obsphys-assembly} (\texttt{model/PppModel.cpp}).  It is
$\sim$1--2\,cm for a typical MEO geometry.

% ============================================================
\subsection{Sagnac (Earth-rotation) correction}\label{sec:obsphys-sagnac}
% ============================================================

The pseudorange relates the satellite position at \emph{signal emission}
to the receiver position at \emph{signal reception}, and both must be
expressed in one ECEF frame.  The ECEF frame rotates with the Earth
during the $\sim$70\,ms transit, so the emission-time satellite
coordinates are rotated into the reception-time frame.  The engine
applies the full rotation, not a scalar approximation.  With the
geometric time-of-flight
$\tau = \lVert\satpos_{\mathrm{em}}-\rxpos_\rx\rVert/\sLight$ and the
Earth rotation rate $\omega_\oplus = 7.2921151467\times10^{-5}$\,rad/s
(\texttt{common/Constants.hpp}), the emission-time position
$\satpos_{\mathrm{em}}$ is rotated by $\theta=\omega_\oplus\tau$ about the
$Z$ axis,
\begin{empheq}[box=\fbox]{equation}
\satpos
=
\bm{R}_z(\omega_\oplus\tau)\,\satpos_{\mathrm{em}},
\qquad
\bm{R}_z(\theta)
=
\begin{bmatrix}
\;\;\cos\theta & \sin\theta & 0\\
-\sin\theta & \cos\theta & 0\\
0 & 0 & 1
\end{bmatrix},
\label{eq:obsphys-sagnac}
\end{empheq}
after which $\geomrange^\sat_\rx=\lVert\satpos-\rxpos_\rx\rVert$ and the
line of sight are computed from the rotated $\satpos$
(\texttt{model/ModelGeometry.cpp}, the \texttt{rotateAndComputeTopo}
lambda).  The Sagnac correction is therefore \emph{part of the geometric
range}, which is why \eqref{eq:obsphys-assembly} has no separate Sagnac
term.  An equivalent scalar form,
$(\omega_\oplus/\sLight)\,(x_\sat y_\rx - y_\sat x_\rx)$, exists in the
code as a \emph{diagnostic} only
(\texttt{GnssCorrections::computeSagnacCorrection},
\texttt{corrections/GnssCorrections.cpp}).  It is the first-order
expansion of the rotation and is not used for modelling.  The
\dcref{sec:products}{Orbit and Clock Product Evaluation document} refers
to this section for the Sagnac rotation.  The same rotation
$\bm{R}_z(\omega(t_{\mathrm{node}}-t_{\mathrm{query}}))$ is used at SP3
interpolation time to rotate ephemeris nodes into the query-epoch frame
(\texttt{model/ModelGeometry.cpp}), for numerical conditioning, as
described with the products.

% ============================================================
\subsection{Group delays and inter-frequency biases}\label{sec:obsphys-bias}
% ============================================================

Two products supply deterministic values for the bias terms of the
Foundations observation equations: the satellite code and phase biases
$B^\sat_\sigl,\,b^\sat_\sigl$ and the receiver biases
$B_{\rx,\sigl},\,b_{\rx,\sigl}$.  This section relates the products to
those symbols.

\paragraph{Broadcast group delay.}
The broadcast navigation message carries a per-system inter-frequency
group delay (GPS TGD, Galileo BGD, BeiDou TGD).  It refers each signal's
code hardware delay to the ionosphere-free combination used to generate
the broadcast clock.  In the uncombined parameterisation these delays
are part of the deterministic code bias applied through the bias term of
\texttt{model/PppModel.cpp}.  They realise $B^\sat_\sigl$ at broadcast
accuracy and are used when no precise bias product is available.

\paragraph{Precise OSB / DCB biases.}
When a precise bias product is loaded, the deterministic correction is the
observable-specific bias (OSB) / differential code bias (DCB) entry for the
exact (\sat,\sigl) pair, time-windowed against its validity interval and
returned by \texttt{GnssCorrections::getBiasCorrection}
(\texttt{corrections/GnssCorrections.cpp}).  This single scalar realises
$B^\sat_\sigl$ (code) or $b^\sat_\sigl$ (phase) for that signal:
\begin{equation}
B_{\mathrm{grp}} = \mathrm{bias}(\sat,\sigl,t),
\qquad
t_{\mathrm{start}} \le t \le t_{\mathrm{end}},
\label{eq:obsphys-bias-window}
\end{equation}
with the application gated by
\texttt{enable\_code\_phase\_bias} (\texttt{model/PppModel.cpp}).  The
bias index (\texttt{GnssCorrections::buildBiasIndex}) excludes two kinds
of entry.

\emph{Zero-valued} code entries, such as the MGEX placeholder
$\mathtt{C2I}/\mathtt{C6I}$ for BeiDou, are treated as \emph{absent}.
Otherwise the model would apply a null correction and also suppress the
estimated bias state for that (\sat,\sigl), leaving a real delay of tens
of metres unmodelled.

\emph{GLONASS code} entries are also treated as missing.  The FDMA
per-channel receiver delay cannot be represented by one value per
(\sat,\sigl), so those signals use the estimated-state path
(\texttt{corrections/GnssCorrections.cpp}).

\paragraph{Estimated vs.\ deterministic biases.}
The estimator handles whatever the deterministic bias does not correct:
mainly the receiver hardware delay $B_{\rx,\sigl},\,b_{\rx,\sigl}$, and
any satellite delay that no product covers.  The receiver biases are
common to all satellites on a signal.  They are absorbed by the receiver
clock and inter-system bias states, or cancel under between-satellite
differencing (\dcref{sec:diff}{the differencing section}).  The
per-(\sat,\sigl) satellite code bias may be \emph{estimated} as a state
when no product
covers it (\texttt{estimation.estimate\_sat\_biases},
\texttt{measurement/MeasurementModel.cpp}).  The state-allocation gate
that prevents estimating a bias the products already supply,
$\mathrm{hasBIA}(\sat,\sigl)$ from the static product, OR'd with
$\mathrm{hasSsrCB}(\sat,\sigl)$ from the live stream, is documented with
the SSR overlay (\dcref{sec:ssr}{the SSR corrections document},
\Cref{eq:ssr-satbias-gate} there); the present document supplies only the
\emph{deterministic} half of that OR.

\paragraph{Relation to the SSR overlay.}
The streamed SSR code/phase OSB (\dcref{sec:ssr}{the SSR corrections
document}) refines the \emph{same} $B^\sat_\sigl,\,b^\sat_\sigl$ terms in
real time, summed onto the predicted observable as \texttt{ssr\_bias\_term}
in \texttt{measurement/MeasurementModel.cpp} alongside the deterministic
$B_{\mathrm{grp}}$ here.  The SSR bias is a per-(\sat,\sigl) scalar like
\eqref{eq:obsphys-bias-window}.  The SSR \emph{orbit} delta is different:
it is expressed in a radial/along/cross (RAC) frame
(\dcref{sec:ssr}{the SSR document}'s \Cref{eq:ssr-rac-frame}), which is
\emph{not} the satellite body frame \eqref{eq:obsphys-bodyframe-matrix}
of this document.  The RAC frame is built from the broadcast position
and velocity to express an orbit correction.  The body frame is built
from the Sun direction to orient the antenna.  The two frames differ in
construction and in purpose.

% ============================================================
\subsection{Provenance and assembly}\label{sec:obsphys-provenance}
% ============================================================

All the deterministic corrections are combined in
\texttt{model/PppModel.cpp}.  \texttt{computeBaseModeledValue}, shared by
the iono-free-pair and uncombined paths, returns the sum
\eqref{eq:obsphys-assembly}.  It calls
\texttt{model/ModelGeometry.cpp} for the Sagnac-rotated geometric range
\eqref{eq:obsphys-sagnac}, \texttt{model/Antenna.cpp} for the total
antenna correction \eqref{eq:obsphys-antenna-total} (which in turn calls
\texttt{model/SatAttitude.cpp} for the body frame and
\texttt{corrections/GnssCorrections.cpp} for the wind-up), and
\texttt{corrections/GnssCorrections.cpp} for the relativistic
\eqref{eq:obsphys-rel}, Shapiro \eqref{eq:obsphys-shapiro} and bias
\eqref{eq:obsphys-bias-window} terms.  The Sun and Moon ephemeris feeding
the attitude frame \eqref{eq:obsphys-yaw-frame}, the Sun-elevation
\eqref{eq:obsphys-beta}, the eclipse predicate
\eqref{eq:obsphys-eclipse} and the wind-up
\eqref{eq:obsphys-windup-dipoles} is produced by
\texttt{model/Planetary.cpp} (an analytic RTKLIB/Meeus series rotated to
ECEF through the IAU 1976/1980 precession-nutation chain, cached per epoch)
and is documented in
\dcref{sec:frames}{the Coordinate Frames and Time Systems document}.  Every
term is a fixed physical model.  The user-facing knobs are enable toggles
and the eclipse-shadow scale; they are cited rather than reproduced because
they set \emph{whether} a correction is applied, never its functional form.

\ifSubfilesClassLoaded{\printbibliography}{}
\end{document}
